\section{Delambre--Gauss analogies, Napier analogies, and the tangent law}
\label{sec:analogies}

The four Delambre formulas are best stated first without division.  This form
remains meaningful even when a traditional quotient presentation would have a
zero denominator.

\begin{theorem}[Delambre--Gauss analogies]
\label{thm:delambre}
For a proper spherical triangle,
\begin{align}
 \sin\frac{A+B}{2}\cos\frac c2
 &=\cos\frac{a-b}{2}\cos\frac C2,                           \tag{7.1}\\
 \sin\frac{A-B}{2}\sin\frac c2
 &=\sin\frac{a-b}{2}\cos\frac C2,                           \tag{7.2}\\
 \cos\frac{A+B}{2}\cos\frac c2
 &=\cos\frac{a+b}{2}\sin\frac C2,                           \tag{7.3}\\
 \cos\frac{A-B}{2}\sin\frac c2
 &=\sin\frac{a+b}{2}\sin\frac C2.                           \tag{7.4}
\end{align}
Cyclic permutation gives twelve cross-multiplied identities.  Whenever the
displayed divisors are apart from zero, these are equivalent to the usual
four quotient analogies.
\end{theorem}

\begin{proof}
We prove (7.1) in detail; the other three use the same calculation with the
corresponding signs.  Squaring and applying the half-angle identities reduces
(7.1) to
\begin{equation}
 (1-\cos(A+B))(1+\cos c)
 =(1+\cos(a-b))(1+\cos C).                                 \tag{7.5}
\end{equation}
The angle cosine law gives
\begin{align*}
 1-\cos(A+B)
 &=1-\cos A\cos B+\sin A\sin B\\
 &=1+\cos C+\sin A\sin B(1-\cos c).
\end{align*}
The side cosine law gives, dually,
\begin{align*}
 1+\cos(a-b)
 &=1+\cos a\cos b+\sin a\sin b\\
 &=1+\cos c+\sin a\sin b(1-\cos C).
\end{align*}
After multiplication by \(1+\cos c\) and \(1+\cos C\), respectively, the
constant terms agree.  The remaining terms agree because
\begin{equation}
 \sin A\sin B\sin^2c=\sin a\sin b\sin^2C,                 \tag{7.6}
\end{equation}
which is the spherical sine law twice.  This proves equality of the squares.
Both sides of (7.1) are nonnegative: \(\cos(c/2),\cos(C/2)>0\),
\(\cos((a-b)/2)>0\) by the strict triangle inequalities, and
\(\sin((A+B)/2)>0\).  Nonnegative-root uniqueness therefore yields (7.1).

For (7.2)--(7.4), square the two sides, replace each half-square by
\((1\pm\cos)/2\), and use the same two cosine laws plus (7.6).  The required
signs follow from the proper-triangle inequalities; signed versions may also
be retained directly when no order simplification is desired.  Every squared
calculation is polynomial after clearing the positive factors.
\end{proof}

The quotient forms most commonly printed are
\begin{align}
 \frac{\sin((A+B)/2)}{\cos(C/2)}
 &=\frac{\cos((a-b)/2)}{\cos(c/2)},                         \tag{7.7}\\
 \frac{\sin((A-B)/2)}{\cos(C/2)}
 &=\frac{\sin((a-b)/2)}{\sin(c/2)},                         \tag{7.8}\\
 \frac{\cos((A+B)/2)}{\sin(C/2)}
 &=\frac{\cos((a+b)/2)}{\cos(c/2)},                         \tag{7.9}\\
 \frac{\cos((A-B)/2)}{\sin(C/2)}
 &=\frac{\sin((a+b)/2)}{\sin(c/2)}.                         \tag{7.10}
\end{align}
Here the half-side and half-angle sines/cosines are positive except for the
possible factors involving \(A-B\) or \(a-b\); those occur only in numerators.

\subsection{Napier's analogies}

\begin{theorem}[Napier analogies]
\label{thm:napier}
With the usual apartness guards,
\begin{align}
 \tan\frac{A+B}{2}
 &=\frac{\cos((a-b)/2)}{\cos((a+b)/2)}\cot\frac C2,         \tag{7.11}\\
 \tan\frac{A-B}{2}
 &=\frac{\sin((a-b)/2)}{\sin((a+b)/2)}\cot\frac C2,         \tag{7.12}\\
 \tan\frac{a+b}{2}
 &=\frac{\cos((A-B)/2)}{\cos((A+B)/2)}\tan\frac c2,        \tag{7.13}\\
 \tan\frac{a-b}{2}
 &=\frac{\sin((A-B)/2)}{\sin((A+B)/2)}\tan\frac c2.        \tag{7.14}
\end{align}
The cyclic permutations give twelve formulas.
\end{theorem}

\begin{proof}
Divide (7.1) by (7.3) to obtain (7.11), and divide (7.2) by (7.4) to obtain
(7.12).  Dividing (7.4) by (7.3) gives (7.13); dividing (7.2) by (7.1) gives
(7.14).  Constructively each division is justified by an explicit apartness
witness.  Alternatively, retain the corresponding cross-multiplied identity,
which follows without any quotient.
\end{proof}

\subsection{The spherical law of tangents}

\begin{corollary}[Law of tangents]
\label{cor:tangent-law}
For each cyclic pair,
\begin{equation}
 \frac{\tan((A-B)/2)}{\tan((A+B)/2)}
 =\frac{\tan((a-b)/2)}{\tan((a+b)/2)},                     \tag{7.15}
\end{equation}
with the required apartness assumptions.  Thus also
\begin{align}
 \frac{\tan((B-C)/2)}{\tan((B+C)/2)}
 &=\frac{\tan((b-c)/2)}{\tan((b+c)/2)},                    \tag{7.16}\\
 \frac{\tan((C-A)/2)}{\tan((C+A)/2)}
 &=\frac{\tan((c-a)/2)}{\tan((c+a)/2)}.                    \tag{7.17}
\end{align}
\end{corollary}

\begin{proof}
Divide (7.12) by (7.11), or equivalently divide (7.14) by (7.13).  The common
factor \(\cot(C/2)\), respectively \(\tan(c/2)\), cancels with its positive
apartness witness.  After cross multiplication, (7.15) is again a finite
polynomial consequence of Theorem~\ref{thm:delambre}.
\end{proof}

\begin{remark}[Constructive preference for cross forms]
The quotient identities are convenient for numerical solution of a triangle,
but their logically strongest versions are the cross-multiplied equations.
They remain valid at symmetric configurations such as \(A=B\) or \(a=b\),
where a printed ratio may become \(0/0\).  The formalization therefore proves
cross forms first and derives quotients only under apartness.
\end{remark}
